\chapter{Theoretical background}
\label{chap:theory}

This chapter gathers the mathematical machinery that the three
architectures of \cref{chap:archi} rest on. The goal is to make the
report \emph{self-contained}: a reader without prior exposure to
deep learning should, after this chapter, be able to read off the
operations performed inside each of our three networks and understand
why they were chosen.

We start with the abstract framing of the task as a windowed
classification problem (\cref{sec:framing}). We then describe in detail
the three building blocks that compose our architectures:
convolutional layers and the supporting normalisations
(\cref{sec:cnn-theory}), residual connections (\cref{sec:resnet-theory}),
and the multi-head self-attention mechanism that powers Transformer
encoders (\cref{sec:attention-theory}). Finally, we summarise the
training ingredients that all three networks share
(\cref{sec:train-theory}).

\section{Framing: time-series classification by windowing}
\label{sec:framing}

Let $x \in \R^N$ denote the gamma-dose signal sampled at uniform
intervals (one minute, in our case). We want to assign each sample
$x_t$ to one of $K$ classes: in the first iteration of the project
$K = 7$ (six anomaly templates plus a \texttt{Normal Background}
class), and in the second iteration the template family was pruned
to $K = 4$.

\paragraph{Windowing.} Operating on a single sample is hopeless: the
classes are defined by the temporal \emph{shape} of an event, not by
the instantaneous value. We therefore process the signal through a
sliding window of length $W$ samples:
\begin{equation}
    \mathbf{x}^{(t)} = (x_{t-W+1}, x_{t-W+2}, \dots, x_t) \in \R^W,
    \label{eq:window}
\end{equation}
and we ask each network to output, from the window $\mathbf{x}^{(t)}$,
a probability distribution
$\hat{p}^{(t)} \in \Delta^{K-1}$ over the $K$ classes\footnote{Here
$\Delta^{K-1}$ denotes the standard $(K-1)$-simplex, i.e.\ the set of
non-negative vectors that sum to $1$.}. In our experiments $W=512$
samples (about $8.5$ hours at $1$ sample / minute).

Successive windows overlap (consecutive predictions share most of their
input), which lets us \emph{accumulate} probabilities at inference time
into a continuous, sample-level confidence trace
(\cref{sec:event-eval}). At training time, windows are extracted at a
larger stride (32 samples) to reduce redundancy.

\paragraph{Labelling rule.}
Because a single window may contain a mixture of background and event
samples, we use an \emph{event-coverage} rule: for every event $e$
overlapping the window, we compute the fraction of $e$'s own samples
that the window captures, and the window inherits the class of the
event with the highest such fraction, provided that fraction exceeds
a threshold $\tau_{\text{cov}}$ (we use $\tau_{\text{cov}} = 0.30$).
Otherwise the window is labelled \texttt{Normal Background}
(see \cref{fig:sliding-window} in \cref{sec:archi-train}). This makes
the boundaries of events slightly fuzzy --- a deliberate choice that
helps the network generalise to events whose exact start and end are
uncertain --- while ensuring that short events (whose total duration
is small compared to $W$) are still seen as positive examples by the
network.

A natural alternative is the \emph{window-coverage} rule: label a
window positive when more than some fraction of \emph{its own}
samples are anomalous. With $W = 512$ and a threshold of $10\%$, this
rule silently maps any event shorter than $\sim 0.1\,W \approx 50$
samples to \texttt{Normal Background} regardless of how completely a
window captures it, because such an event simply cannot fill enough
of the window to clear the threshold. The v2 generator deliberately
produces events as short as $\sim 30$ samples, so anchoring the
threshold to the event's own length rather than the window's avoids
losing them at the labelling stage.

\paragraph{Mathematical objective.}
Letting $f_\theta : \R^W \to \R^K$ denote the network with parameters
$\theta$, the empirical objective is the (weighted) cross-entropy
\begin{equation}
    \mathcal{L}(\theta) = - \frac{1}{|\mathcal{D}|} \sum_{(\mathbf{x}, y) \in \mathcal{D}}
        w_y \log \softmax(f_\theta(\mathbf{x}))_y,
    \label{eq:cross-entropy}
\end{equation}
where $\mathcal{D}$ is the training set of (window, label) pairs,
$y \in \{0, \dots, K-1\}$ the integer label, and $w_y > 0$ a class
weight that compensates for class imbalance, defined in
\cref{sec:cw-theory}.

\section{Convolutional layers and supporting normalisations}
\label{sec:cnn-theory}

The cornerstone of all three architectures is the
\textbf{one-dimensional convolution}.

\subsection{The 1D convolution operator}
\label{sec:conv-theory}

Let $\mathbf{u} \in \R^{C_{\text{in}} \times L}$ denote an input
sequence with $C_{\text{in}}$ channels and length $L$, and let
$\mathbf{w} \in \R^{C_{\text{out}} \times C_{\text{in}} \times k}$
denote a learnable kernel of $C_{\text{out}}$ output channels with
support of size $k$. The convolution output
$\mathbf{v} \in \R^{C_{\text{out}} \times L'}$ has entries:
\begin{equation}
    \mathbf{v}_{c, t} = \sum_{c'=1}^{C_{\text{in}}} \sum_{i=1}^{k}
        \mathbf{w}_{c, c', i} \cdot \mathbf{u}_{c', \, s \cdot t + i - p} + b_c,
    \label{eq:conv}
\end{equation}
where $b_c$ is a learnable bias, $s$ is the \emph{stride} and $p$ is the
\emph{padding}. The output length is
$L' = \lfloor (L + 2p - k)/s \rfloor + 1$. Three observations matter:
\begin{itemize}[leftmargin=*, itemsep=0.2em]
    \item The kernel $\mathbf{w}_{c, \cdot, \cdot}$ is shared across all
    time positions $t$, so the operator is \emph{translation
    equivariant}: shifting the input shifts the output by the same
    amount. This is exactly what we want for an anomaly that may occur
    anywhere in the window.
    \item The number of parameters is $C_{\text{out}} \cdot C_{\text{in}} \cdot k$
    (plus $C_{\text{out}}$ biases), \emph{independent of} $L$. Convolutions
    are therefore much cheaper than the dense fully-connected
    alternative, which would have $C_{\text{out}} \cdot C_{\text{in}} \cdot L \cdot L'$
    parameters.
    \item Each output channel $c$ acts as a learned \emph{feature
    detector}: it lights up where the input matches a particular local
    pattern of length $k$.
\end{itemize}

\paragraph{Receptive field.}
Stacking several convolutional layers, optionally with sub-sampling
(stride or pooling, see below), the \emph{receptive field} of a deeper
neuron --- the range of input samples that can affect its value ---
grows roughly linearly with depth and exponentially with the
sub-sampling factor. Our $W=512$ window is short enough that three
stride-2 pooling steps (giving a receptive field of $\sim 50$ samples on
top of $k=7$ kernels) already cover the duration of the smallest
expected anomalies (period $\geq 200$ samples once temporal pooling is
accounted for).

\subsection{Pointwise nonlinearity: ReLU}
\label{sec:relu-theory}

Stacking purely linear convolutions yields, by composition, another
linear operator and so cannot learn rich features. Between successive
convolutions we therefore apply an element-wise non-linearity. We use
the \emph{rectified linear unit}, $\ReLU(x) = \max(0, x)$, which is
cheap to compute, has a non-vanishing gradient on its positive half-line,
and has become the de facto standard activation in convolutional
networks. The two normalisation variants used below already absorb
the negative-half collapse problem (``dying ReLU''), so we do not need a
fancier activation.

\subsection{Pooling}
\label{sec:pool-theory}

After a few convolutions we typically reduce the temporal resolution by
\emph{pooling} a contiguous group of values into one. Our networks use
two kinds of pooling.

\paragraph{Max-pooling.} Of the values in a window of size $k_p$,
keep only the maximum:
$\mathrm{MaxPool}_k(u)_t = \max_{i \in [1, k_p]} u_{s_p t + i}$,
with stride $s_p = k_p$ in our case so that successive max-pools do not
overlap. The output is shorter by a factor $k_p$. Max-pooling makes
the representation locally invariant to small temporal shifts
of an anomaly.

\paragraph{Adaptive average pooling.} Used near the end of the network
to bring the temporal axis down to a fixed length regardless of the
input window size: it averages the input into a target number of bins.
Our networks emit $4$ temporal bins from \texttt{AdaptiveAvgPool1d(4)},
which preserves a coarse temporal shape (``did the activation peak in
the first or the last quarter of the window?'') while removing the
exact location, before the classifier head.

\subsection{Batch normalisation}
\label{sec:bn-theory}

A common pathology of deep networks is that the distribution of
activations inside the network shifts during training, slowing down
convergence (the so-called \emph{internal covariate shift}). Batch
normalisation \cite{ioffe2015batch} addresses this by standardising
each channel of each layer over the current mini-batch:
\begin{equation}
    \BN(\mathbf{u})_{c, t} = \gamma_c \cdot \frac{\mathbf{u}_{c, t} - \mu_c}{\sqrt{\sigma_c^2 + \varepsilon}} + \beta_c,
    \label{eq:bn}
\end{equation}
where $\mu_c, \sigma_c^2$ are the channel mean and variance computed
over the batch dimension (and the time dimension for 1D-BN), and
$\gamma_c, \beta_c$ are learnable rescaling parameters.

At inference time, $\mu_c$ and $\sigma_c^2$ are replaced by exponential
moving averages collected during training. This means BN behaves
\emph{differently} at training and inference time, which can become
problematic when there is a distribution shift between training and
deployment --- a subtlety that motivates the GroupNorm variant in our
two more advanced architectures.

\subsection{Group normalisation}
\label{sec:gn-theory}

Group normalisation \cite{wu2018group} performs the same affine
standardisation \emph{but} computes the statistics
$\mu, \sigma^2$ over a partition of the channels into $G$ groups,
independently for each sample. There are no running statistics to
maintain; train and inference behaviour are identical, and the
operator becomes amplitude-invariant in a window-by-window fashion.
This robustness against shifts in absolute scale is exactly the
property we need in a stem that processes signals with a fluctuating
baseline. Our ResNet 1D and CNN+Attention models use GroupNorm with $8$
groups (or fewer for narrow channel widths) throughout their
convolutional stems.

\section{Residual connections}
\label{sec:resnet-theory}

\paragraph{The vanishing-gradient problem.}
A pure stack of convolutional layers $f_1 \circ f_2 \circ \dots \circ f_L$
becomes harder to optimise as $L$ grows: the gradient of the loss with
respect to early-layer parameters is a product of $L$ Jacobian factors,
and unless each of them is exactly orthogonal it tends to either
vanish or explode exponentially.

\paragraph{Residual blocks.}
He et al.\ \cite{he2016deep} proposed to reformulate a stack of two or
three layers as a learnable \emph{residual} to the identity:
\begin{equation}
    \mathrm{ResBlock}(\mathbf{u}) = \ReLU\bigl(\, \mathbf{u} + \mathcal{F}(\mathbf{u}) \,\bigr),
    \label{eq:resblock}
\end{equation}
where $\mathcal{F}$ is a small sub-network (here: $\Conv \to \GN \to
\ReLU \to \Conv \to \GN$). The shortcut $+\mathbf{u}$ guarantees that
the block can degenerate to the identity ($\mathcal{F}=0$) at no cost,
which means adding a residual block to a working network cannot make
things worse. In gradient terms the backward pass propagates
\emph{additively} through the shortcut, so deep stacks remain
trainable.

When the input and the output of $\mathcal{F}$ have different channel
counts, the shortcut is replaced by a $1{\times}1$ convolution that
projects them into compatible spaces. We use this only at block
transitions, where the channel width grows from $16 \to 32 \to 64$.

\paragraph{Why residual learning helps anomaly classification.}
For our problem the residual block can be intuitively read as ``copy
the input through, then \emph{add} a learnt correction'': in early
blocks the correction extracts local shape primitives (peaks, ramps);
in later blocks it refines them by re-using the original signal context
that has been preserved by the shortcut.

\section{Attention and the Transformer encoder}
\label{sec:attention-theory}

The third architecture in this report augments a small convolutional
stem with a \emph{Transformer encoder} \cite{vaswani2017attention}.
Transformers are now the dominant architecture in natural language and
sequence modelling; they have also proved very competitive on
time-series classification \cite{hu2022eeg}. The mechanism that powers
them is \emph{self-attention}, which we now describe carefully.

\subsection{Self-attention}
\label{sec:selfattn-theory}

The input to a self-attention layer is a sequence of $T$ token vectors
$\mathbf{Z} \in \R^{T \times d}$ (one vector per time step, of dimension
$d$). The layer first projects $\mathbf{Z}$ into three learnable
representations called \emph{queries}, \emph{keys} and \emph{values}:
\begin{equation}
    \mathbf{Q} = \mathbf{Z} \mathbf{W}^Q,
    \quad
    \mathbf{K} = \mathbf{Z} \mathbf{W}^K,
    \quad
    \mathbf{V} = \mathbf{Z} \mathbf{W}^V,
    \label{eq:qkv}
\end{equation}
where $\mathbf{W}^Q, \mathbf{W}^K, \mathbf{W}^V \in \R^{d \times d_h}$
are learnable parameter matrices ($d_h$ is the per-head dimension).

The attention output is then
\begin{equation}
    \Attn(\mathbf{Q}, \mathbf{K}, \mathbf{V})
    = \softmax\!\left( \frac{\mathbf{Q} \mathbf{K}^\top}{\sqrt{d_h}} \right) \mathbf{V}.
    \label{eq:attention}
\end{equation}
The matrix $\mathbf{A} = \softmax(\mathbf{Q} \mathbf{K}^\top / \sqrt{d_h})
\in \R^{T \times T}$ has rows that sum to $1$; its entry $A_{ij}$ is
the fraction of token $i$'s output that comes from token $j$. The
output sequence at row $i$ is therefore a \emph{convex combination of
all input value vectors}, weighted by how compatible the query of token
$i$ is with the key of each other token. The scaling by $\sqrt{d_h}$
keeps the dot products from saturating the softmax in high dimensions.

\paragraph{Why this matters for our problem.}
Convolutional layers have a bounded receptive field: a neuron in layer
$\ell$ depends only on samples within a window of size proportional to
$\ell$. Self-attention, by contrast, lets every output position attend
to \emph{every} input position in a single layer, with a
\emph{learnt}, data-dependent weighting. This long-range, content-based
mixing is exactly what is needed to disambiguate, say, a sharp rise
that could be a \texttt{fast\_ascend} or just a noise excursion, by
looking at what happens \emph{after} the rise within the same window.

\subsection{Multi-head attention}
\label{sec:mha-theory}

A single attention head can only learn one notion of similarity. To
let the model attend to several relations in parallel, multi-head
attention runs $h$ independent attention heads on parallel projections
of the same input, and concatenates their outputs:
\begin{equation}
    \MHA(\mathbf{Z}) = \mathrm{Concat}(\mathbf{O}_1, \dots, \mathbf{O}_h) \mathbf{W}^O,
    \qquad
    \mathbf{O}_j = \Attn(\mathbf{Z} \mathbf{W}_j^Q,\, \mathbf{Z} \mathbf{W}_j^K,\, \mathbf{Z} \mathbf{W}_j^V).
    \label{eq:mha}
\end{equation}
We use $h=4$ heads with $d_h = d / h = 16$, so the total number of
attention parameters is the same as that of a single head with
dimension $d=64$.

\subsection{Positional encoding}
\label{sec:posenc-theory}

The attention operation \cref{eq:attention} is \emph{permutation
equivariant} in the input: shuffling the tokens of $\mathbf{Z}$
shuffles the output rows in the same way. For time series this is a
problem: we explicitly want the network to know whether a feature
occurred at the start or at the end of the window.

We follow Vaswani et al.\ \cite{vaswani2017attention} and add a fixed
\emph{sinusoidal positional encoding} to the token sequence before
feeding it to the encoder:
\begin{equation}
    \mathrm{PE}_{t, 2k}   = \sin\!\left( t \cdot \omega^{-2k/d} \right),
    \quad
    \mathrm{PE}_{t, 2k+1} = \cos\!\left( t \cdot \omega^{-2k/d} \right),
    \quad
    \omega = 10\,000.
    \label{eq:posenc}
\end{equation}
This vector is added to the token at position $t$ before any attention
layer is applied. Sines and cosines at geometrically decreasing
frequencies let the network express both fine and coarse positional
relations as easily-learned linear combinations.

\subsection{The Transformer encoder block}
\label{sec:encoder-theory}

The complete \emph{Transformer encoder layer} composes multi-head
self-attention with a two-layer feed-forward network and two LayerNorm
\cite{ba2016layer} steps wrapped in residual connections
(\cref{fig:transformer-block}):
\begin{align}
    \mathbf{Z}' &= \mathbf{Z} + \MHA(\mathrm{LN}(\mathbf{Z})), \\
    \mathbf{Z}'' &= \mathbf{Z}' + \mathrm{FFN}(\mathrm{LN}(\mathbf{Z}')),
\end{align}
with $\mathrm{FFN}(x) = \mathbf{W}_2 \, \mathrm{GELU}(\mathbf{W}_1 x + \mathbf{b}_1) + \mathbf{b}_2$.
The pre-LayerNorm convention used here (LN \emph{before} attention/FFN)
is more stable to train than the original post-LN convention.

\paragraph{The CLS aggregation token.}
After the encoder we still have a length-$T$ sequence, but we need a
single class probability. Following BERT \cite{devlin2019bert}, we
prepend to the sequence a special learnable vector --- the
\emph{classification token} \texttt{[CLS]} --- which has no signal
content of its own but is allowed to attend to every other token. After
the encoder, the output row at the \texttt{[CLS]} position has
aggregated information from the entire sequence and is fed alone to the
linear classifier.

\begin{figure}[H]
    \centering
    \begin{tikzpicture}[
        font=\small,
        block/.style={draw, rounded corners=2pt, align=center, minimum height=0.7cm, inner sep=4pt},
        attn/.style={block, fill=red!12},
        ffn/.style={block, fill=blue!12},
        norm/.style={block, fill=black!8},
        ar/.style={-{Stealth[length=2mm]}}
    ]
        \node[norm] (ln1) {LayerNorm};
        \node[attn, right=0.7cm of ln1] (mha) {Multi-Head\\Self-Attention};
        \node[circle, draw, right=0.7cm of mha, minimum size=4mm, inner sep=0pt] (add1) {$+$};
        \node[norm, right=0.6cm of add1] (ln2) {LayerNorm};
        \node[ffn, right=0.7cm of ln2] (ffn) {Feed-Forward\\(2-layer GELU)};
        \node[circle, draw, right=0.7cm of ffn, minimum size=4mm, inner sep=0pt] (add2) {$+$};

        \node[left=0.7cm of ln1] (in) {$\mathbf{Z}$};
        \node[right=0.7cm of add2] (out) {$\mathbf{Z}''$};

        \draw[ar] (in) -- (ln1);
        \draw[ar] (ln1) -- (mha);
        \draw[ar] (mha) -- (add1);
        \draw[ar] (add1) -- (ln2);
        \draw[ar] (ln2) -- (ffn);
        \draw[ar] (ffn) -- (add2);
        \draw[ar] (add2) -- (out);

        % residual 1
        \draw[ar] (in) to[out=-30, in=-90] ($(add1.south)$);
        % residual 2
        \coordinate (rmid) at ($(add1)!0.5!(ln2)$);
        \draw[ar] (rmid) to[out=-30, in=-90] ($(add2.south)$);
    \end{tikzpicture}
    \caption{A pre-LN Transformer encoder block. Each sub-layer is wrapped in a residual connection; LayerNorm is applied before the sub-layer (pre-norm).}
    \label{fig:transformer-block}
\end{figure}

\section{Training fundamentals}
\label{sec:train-theory}

All three networks are trained with the same recipe.

\subsection{Cross-entropy loss}
\label{sec:ce-theory}

For multi-class classification, the standard loss is the
\emph{categorical cross-entropy} introduced in \cref{eq:cross-entropy}.
Letting $\mathbf{z} = f_\theta(\mathbf{x}) \in \R^K$ be the network's
output logits and $y$ the integer label, the contribution of a single
example to the loss is
\begin{equation}
    \mathcal{L}(\mathbf{x}, y) = -\log \frac{\exp(z_y)}{\sum_{c=1}^{K} \exp(z_c)},
\end{equation}
i.e.\ minus the logarithm of the probability the network assigns to the
correct class after a softmax normalisation. The gradient of $\mathcal{L}$
with respect to $z_c$ is $\softmax(\mathbf{z})_c - \mathbf{1}[c=y]$:
the network is encouraged to push the correct logit up and pull the
others down, in proportion to how confident the current (wrong) belief
is.

\subsection{Class weighting}
\label{sec:cw-theory}

When the training set is imbalanced, the loss is dominated by the
majority class and minority classes are poorly learned. We compensate
with a per-class weight (a per-example multiplier on the loss)
inversely proportional to the class frequency,
\begin{equation}
    w_c = \frac{N}{K \cdot n_c},
\end{equation}
where $n_c$ is the count of class $c$ in the training set and $N$ is
the total count. The numerical values that this rule produces are
reported, for both v1 and v2, in the dedicated training sections; they
will play a non-trivial role at deployment because the
training-vs-deployment class prior shift is the dominant cause of v1's
failure on real data (\cref{chap:v1real}).

\subsection{Adam optimiser}
\label{sec:adam-theory}

The model is fit by stochastic gradient descent on mini-batches with
the Adam optimiser \cite{kingma2015adam}, which maintains, in addition
to the gradient, an exponential moving average of the first and second
moments of the gradient per parameter and rescales the update
accordingly:
\begin{equation}
    m_t = \beta_1 m_{t-1} + (1 - \beta_1) g_t,
    \quad
    v_t = \beta_2 v_{t-1} + (1 - \beta_2) g_t^2,
    \quad
    \theta_{t+1} = \theta_t - \eta \, \frac{\hat{m}_t}{\sqrt{\hat{v}_t} + \varepsilon},
\end{equation}
where $g_t$ is the gradient at step $t$ and $\hat{m}_t, \hat{v}_t$ are
the bias-corrected moments. The defaults $\beta_1=0.9, \beta_2=0.999,
\varepsilon=10^{-8}$ are used. The initial learning rate is $\eta=10^{-3}$.

\subsection{Learning rate scheduling}
\label{sec:sched-theory}

A constant learning rate is rarely optimal: large steps early in
training are useful to escape poor local minima, but late in training
they oscillate around the minimum. We use
\texttt{ReduceLROnPlateau}: whenever the validation loss has not
improved for two consecutive epochs, the learning rate is multiplied by
$0.5$. This is the simplest form of adaptive schedule and is enough
here to drive the loss down by roughly $20\%$ in the last third of
training.

\subsection{Early stopping}
\label{sec:es-theory}

To avoid over-fitting we monitor the validation loss and stop training
when it has not improved by more than $\delta=5\cdot 10^{-3}$ for ten
consecutive epochs. The model weights with the lowest validation loss
are then re-loaded.

\subsection{Temporal train/validation/test split}
\label{sec:split-theory}

A subtle pitfall of sliding-window time-series classification is
\emph{data leakage}: if windows are extracted first and then shuffled
across train/test, two windows that overlap in time can end up in
different splits, so the test set effectively includes parts of the
training signal. We avoid this by splitting the \emph{raw signal}
along time first (the first $80\%$ becomes the train+validation pool,
the remaining $20\%$ becomes the test set), and only then applying the
sliding-window extraction independently to each split. The
$80\%$ train+val pool is then further split uniformly at random into
$80\%/20\%$, which is safe because no overlap is possible anymore
between training and validation windows on the same signal slice if
the stride matches.

\subsection{Window normalisation}
\label{sec:norm-theory}

Because the baseline of the signal drifts slowly
(\cref{sec:real-data}), feeding the absolute dose to the network is
inefficient: the network would waste capacity to learn the same
features at every possible baseline. We \emph{centre} each window
independently --- subtract its own mean --- and divide by a single
global standard deviation $\sigma_{\text{train}}$ estimated on the
centred training set:
\begin{equation}
    \tilde{\mathbf{x}}^{(t)} = \frac{\mathbf{x}^{(t)} - \bar{\mathbf{x}}^{(t)}}{\sigma_{\text{train}}}.
    \label{eq:window-norm}
\end{equation}
The same $\sigma_{\text{train}}$ is used at inference; saving it in a
JSON file alongside the model weights guarantees deterministic and
reproducible scaling between training and deployment. This window-level
centring is a strong inductive bias --- the network only ever sees
deviations from the local baseline --- and the consequences of that
bias on long sustained events will surface in \cref{chap:v1real}.
